\documentclass{ctexart}
\usepackage{amsmath, amssymb}
\usepackage{geometry}
\geometry{a4paper, margin=2cm}

\title{\textbf{第6章 刚体力学} --- 例题}
\author{}
\date{}

\begin{document}
\maketitle

\begin{enumerate}

\item（书中例题 5.1，p.182）曲柄连杆机构，曲柄长度为 $r$，与 $x$ 轴的夹角 $\varphi = \omega t$，其中 $\omega$ 为常量。求 $T$ 形连杆在 $t$ 时刻的速度和加速度。

\item 已知棒长 $L$、质量 $M$，在摩擦系数为 $\mu$ 的桌面转动。求摩擦力对 $y$ 轴的力矩。

\item（书中例题 6.1，p.198）
\begin{enumerate}
  \item 求均匀细棒绕端点轴的转动惯量；
  \item 求均匀细棒绕中心轴的转动惯量；
  \item 求圆环绕中心轴的转动惯量；
  \item 求圆盘绕中心轴的转动惯量。
\end{enumerate}

\item 已知书中例题 6.1 求了杆通过中心轴的转动惯量，用平行轴定理求均匀细棒一端点的转动惯量。

\item 求过圆盘边缘与圆盘中心轴平行的轴的转动惯量。

\item（重点）半径为 $R$、长为 $L$、质量为 $M$ 的实心圆柱体对中心直径的转动惯量。

\item（书中例题 6.3，p.202）已知滑轮半径为 $R$、质量为 $M$，绳子不可伸缩，绳子与滑轮间无滑动，轴处无摩擦，两个悬挂物的质量分别为 $m_1$、$m_2$。求两重物的加速度、滑轮的角加速度、绳中的张力。

\item（书中例题 6.4，p.202）两个皮带轮半径分别为 $R_1$、$R_2$，质量分别为 $m_1$、$m_2$，分别绕固定轴 $O_1$、$O_2$ 转动，用皮带相连，轮 $1$ 作用力矩 $M_1$，轮 $2$ 有负载力矩 $M_2$，皮带与轮无滑动，轴处无摩擦。求轮 $1$ 的角加速度。

\item（书中例题 6.5，p.202）飞轮齿轮 $1$ 绕转轴 $1$ 的转动惯量 $J_1 = 98.0\,\text{kg}\cdot\text{m}^{2}$，飞轮齿轮 $2$ 绕转轴 $2$ 的转动惯量 $J_2 = 78.4\,\text{kg}\cdot\text{m}^{2}$，两齿轮咬合传动，齿数比 $Z_1:Z_2 = 3:2$，$r_1 = 10\,\text{cm}$，轴 $1$ 从静止在 $10\,\text{s}$ 匀加速到 $1500\,\text{r/min}$。求加在轴 $1$ 上的力矩 $M$ 和齿轮间的相互作用力 $Q$。

\item 一根长为 $l$、质量为 $m$ 的均匀细直棒，可绕轴 $O$ 在竖直平面内转动，初始时它在水平位置。求它由此下摆 $\theta$ 角时的角速度 $\omega$。

\item 圆盘以 $\omega_0$ 在桌面上转动，受摩擦力而静止。求圆盘静止所需时间。

\item（书上例题 6.6）长为 $l$ 质量分布不均匀的细杆，其线密度为 $\lambda = a + br$（$a$、$b$ 为常量），杆可绕 $A$ 端的 $z$ 轴在铅垂面内转动。将杆从水平位置释放，求杆转到铅垂位置过程中重力做的功。

\item（书中例题 6.8，p.209）圆盘滑轮质量 $M$、半径 $R$，绕轻绳，绳的另一端系一质量 $m$ 的物体，轴无摩擦，开始时系统静止。求物体下降 $s$ 时滑轮的角速度和角加速度。

\item（书中例题 6.13，p.217）长 $l$、质量 $M$，铅直悬挂，初始处于静止状态，杆的中心受一冲量 $I$ 作用，方向与杆垂直。求冲量作用结束时杆的角速度。

\item（书中例题 6.15，p.220）半径为 $R$、质量为 $4m$ 且均匀分布的转盘可绕铅垂轴无摩擦地转动，盘上站有质量为 $m$ 的四个人，两个人站在盘的边沿，两个人站在 $R/2$ 处，整个系统开始以角速度 $\omega_0$ 转动。此后，人相对于盘作圆周运动，站在边沿上的两人相对于盘的速度为 $\mu$，站在 $R/2$ 处的两人相对于盘的速度为 $2\mu$。求：
\begin{enumerate}
  \item 人走动后转盘的角速度；
  \item 要使转盘停止转动，$\mu$ 的大小。
\end{enumerate}

\item（书中例题 6.16，p.221，重点）长为 $L$、质量为 $M$ 的均匀杆，一端悬挂，由水平位置无初速度地下落，在铅直位置与质量为 $m$ 的物体 $A$ 作完全非弹性碰撞，碰后物体 $A$ 沿摩擦系数为 $\mu$ 的水平面滑动。求物体 $A$ 滑动的距离。

\item 一个刚体系统，转动惯量 $J = \dfrac13 ml^{2}$，现有一水平力作用于距轴为 $l'$ 处。求轴对棒的作用力（轴反力）。

\item 一长为 $l$ 的匀质细杆，可绕通过中心的固定水平轴在铅垂面内自由转动，开始时杆静止于水平位置。一质量与杆相同的昆虫以速度 $v_0$ 垂直落到距点 $O$ 为 $l/4$ 处的杆上，昆虫落下后立即向杆的端点爬行，如图所示。若要使杆以匀角速度转动，求昆虫沿杆爬行的速度。

\end{enumerate}

\end{document}
